\begin{table}[htbp]
\centering
\caption{Round-Number Heaping at Non-Regulatory Thresholds}
\label{tab:heaping}
\begin{tabular}{lccc}
\hline\hline
Bed count & Observed & Avg.\ neighbors & Heaping ratio \\
\hline
20 & 1,148 & 576 & 1.99 \\
30 & 754 & 299 & 2.52 \\
40 & 1,202 & 343 & 3.51 \\
60 & 977 & 236 & 4.15 \\
70 & 357 & 225 & 1.59 \\
80 & 481 & 182 & 2.65 \\
90 & 336 & 210 & 1.60 \\
110 & 279 & 216 & 1.29 \\
120 & 256 & 175 & 1.46 \\
\hline
Average & & & 2.31 \\
\hline
\multicolumn{4}{l}{\textit{Regulatory thresholds (for comparison):}} \\
25 (CAH) & 10,264 & 345 & 29.8 \\
50 (RHC/REH) & 838 & 149 & 5.6 \\
100 (DSH) & 419 & 241 & 1.7 \\
\hline\hline
\end{tabular}
\begin{tablenotes}
\small
\item \textit{Notes:} Heaping ratio = count at round number / average count at $\pm 1$ and $\pm 2$ neighbors. Non-regulatory round numbers are multiples of 10 that do not coincide with Medicare payment thresholds. The 25-bed CAH threshold shows a heaping ratio 13$\times$ the non-regulatory average, while 50-bed and 100-bed thresholds are within the normal heaping range.
\end{tablenotes}
\end{table}

